% !TeX program = pdflatex
% Загрузите main.tex, Paper.tex и gost.sty в один проект.
% Основной документ: main.tex. Компилятор: pdfLaTeX.
\documentclass[14pt,a4paper,oneside]{extreport}
\usepackage{gost}

% Заполните свои данные. Текст отчёта редактируется в Paper.tex.
\newcommand{\LabNumber}{2}
\newcommand{\LabTitle}{Сборка образа: Dockerfile, оптимизация, healthcheck и теги}
\newcommand{\CourseName}{Контейнеризация и оркестрация приложений}
\newcommand{\StudentName}{Ивановский Роман Евгеньевич}
\newcommand{\StudentShortName}{Ивановский Р.\,Е.}
\newcommand{\StudentISU}{466021}
\newcommand{\StudentGroup}{К3320}
\newcommand{\FacultyName}{Факультет прикладной информатики}
\newcommand{\DegreeProgramme}{11.03.02 Инфокоммуникационные технологии и системы связи}
\newcommand{\EducationProgramme}{Программирование в инфокоммуникационных системах}
\newcommand{\TeacherName}{Аминов Н.\,С.}
\newcommand{\ReportCity}{Санкт-Петербург}
\newcommand{\ReportYear}{2026}
\newcommand{\RepositoryURL}{https://github.com/st1verka/course-cont-orch-app/tree/lab-2}

\hypersetup{
  pdftitle={Отчёт по лабораторной работе \LabNumber},
  pdfauthor={\StudentName},
  pdfsubject={\CourseName}
}

\begin{document}
\begin{titlepage}
  \thispagestyle{empty}
  \begin{center}
    \fontsize{12pt}{16pt}\selectfont
    Министерство науки и высшего образования Российской Федерации\\[0.25\baselineskip]
    ФЕДЕРАЛЬНОЕ ГОСУДАРСТВЕННОЕ АВТОНОМНОЕ\\
    ОБРАЗОВАТЕЛЬНОЕ УЧРЕЖДЕНИЕ ВЫСШЕГО ОБРАЗОВАНИЯ\\[0.25\baselineskip]
    НАЦИОНАЛЬНЫЙ ИССЛЕДОВАТЕЛЬСКИЙ УНИВЕРСИТЕТ ИТМО\\
    ITMO University
  \end{center}
  \vspace{0.8\baselineskip}
  \begin{center}
    \textbf{ОТЧЁТ ПО ЛАБОРАТОРНОЙ РАБОТЕ №\LabNumber}
  \end{center}
  \vspace{0.4\baselineskip}
  \begingroup
  \setlength{\parindent}{0pt}
  По дисциплине: \CourseName\par\medskip
  Тема работы: \LabTitle\par\medskip
  Обучающийся: \StudentName\par\medskip
  Номер ИСУ: \StudentISU\par\medskip
  Факультет: \FacultyName\par\medskip
  Группа: \StudentGroup\par\medskip
  Направление подготовки: \DegreeProgramme\par\medskip
  Образовательная программа:\\\EducationProgramme\par
  \endgroup
  \vfill
  \begin{tabularx}{\textwidth}{@{}lXc@{}}
    Обучающийся & & \StudentShortName\\
    & & \rule{38mm}{0.4pt}\\[-0.25\baselineskip]
    & & {\fontsize{12pt}{14pt}\selectfont дата, подпись}\\[0.6\baselineskip]
    Руководитель & & \TeacherName\\
    & & \rule{38mm}{0.4pt}\\[-0.25\baselineskip]
    & & {\fontsize{12pt}{14pt}\selectfont дата, подпись}
  \end{tabularx}
  \vfill
  \begin{center}
    \ReportCity\\
    \ReportYear
  \end{center}
\end{titlepage}
\setcounter{page}{2}
\tableofcontents
\reportelement{Введение}
Цель работы --- освоить сборку Docker-образа Python-приложения, сравнить варианты оптимизации, исключить служебные файлы из контекста сборки, настроить проверку здоровья контейнера и тегирование.
В качестве приложения использован персональный ассистент на FastAPI из репозитория \url{https://github.com/FilBTi/assist-py}. Работа выполнена локально в Docker Desktop.

\reportelement{Персональные параметры}
\begin{table}[H]
\caption{Персональные параметры лабораторной работы}
\begin{tabularx}{\textwidth}{|p{48mm}|X|}
\hline
Параметр & Значение \\\hline
ФИО & Ивановский Роман Евгеньевич \\\hline
Группа & К3320 \\\hline
Номер ИСУ & 466021 \\\hline
Персональный порт & $5430+21=5451$ \\\hline
Префикс & \texttt{ivanovsky} \\\hline
Учебный секрет & \texttt{secret21ivanovsky} \\\hline
\end{tabularx}
\end{table}
Секрет используется только для учебного эндпоинта и передаётся контейнеру через переменную окружения APP\_SECRET; в Dockerfile он не записывается.

\chapter{Первый рабочий образ}
\section{Подготовка приложения и Dockerfile}
Создан каталог проекта и клонировано исходное приложение в solutions/assist-py/.
\reportscreenshot{images/step-01.png}{Клонирование исходного приложения}{fig:s01}

В корне приложения создан Dockerfile на основе python:3.11-slim. Указаны метаданные автора и версии, рабочая директория /app и персональный порт 5451. Сначала копируется requirements.txt и устанавливаются зависимости, затем копируется код приложения.
Команда запуска использует shell-подстановку PORT и exec для передачи сигналов процессу uvicorn.
\reportscreenshot{images/step-02.png}{Базовый Dockerfile}{fig:s02}

\section{Сборка и запуск}
Образ собран командой:
\begin{code}
docker build -t ivanovsky-assistant:v1 .
\end{code}
\reportscreenshot{images/step-03.png}{Успешная сборка образа v1}{fig:s03}

Сборка завершилась успешно. Предупреждение JSONArgsRecommended относится к shell-форме CMD; она используется для подстановки переменной PORT, как предусмотрено методичкой.
Контейнер запущен с именем ivanovsky-assistant-1, пробросом 5451:5451 и переменными APP\_OWNER и APP\_SECRET.
\reportscreenshot{images/step-04.png}{Запуск первого контейнера}{fig:s04}

\section{Проверка эндпоинтов}
Главная страница вернула приветствие владельцу ivanovsky и сведения о сервисе.
\reportscreenshot{images/step-05.png}{Ответ главной страницы}{fig:s05}

Эндпоинт /me вернул владельца ivanovsky и префикс учебного секрета. Эта проверка дополнительно выполнена после запуска финального контейнера.
\reportscreenshot{images/me.png}{Персонализация в ответе /me}{fig:me}
Создана заметка «ЛР2: сборка образа». Запрос списка заметок подтвердил её сохранение.
\reportscreenshot{images/step-06.png}{Создание и чтение заметки}{fig:s06}

Запрос с правильным учебным ключом вернул флаг:\par\noindent\texttt{LAB2\_SUCCESS\_IVANOVSKY}.
\reportscreenshot{images/step-07.png}{Проверка секретного эндпоинта}{fig:s07}
\reportscreenshot{images/step-08.png}{Общий вывод проверок первого контейнера}{fig:s08}

\chapter{Оптимизация образа}
\section{Отключение кеша pip}
Создан Dockerfile.optimized, в котором установка зависимостей выполняется с флагом --no-cache-dir. Он запрещает сохранение кеша pip при установке пакетов.
\reportscreenshot{images/step-09.png}{Dockerfile.optimized с отключением кеша pip}{fig:s09}

Оптимизированный образ собран с тегом v2. По столбцу DISK USAGE размер v1 составляет 313 МБ, а v2 --- 287 МБ. По округлённым значениям уменьшение составило 26 МБ, около 8,3\% от размера v1.
\reportscreenshot{images/step-10.png}{Сравнение размеров v1 и v2}{fig:s10}

Файл requirements.txt копируется отдельно. Это позволяет повторно использовать слой установки зависимостей, если изменился только код приложения. При копировании всего проекта перед установкой пакетов изменение бизнес-логики инвалидировало бы этот слой.

\section{Многоэтапная сборка}
На стадии builder создаётся виртуальное окружение /opt/venv и устанавливаются зависимости без кеша pip. На итоговой стадии используется тот же базовый образ, копируются виртуальное окружение и приложение. uvicorn запускается по полному пути /opt/venv/bin/uvicorn.
\reportscreenshot{images/step-11.png}{Dockerfile с двумя стадиями сборки}{fig:s11}

Образ с тегом ms занял 307 МБ, на 20 МБ больше v2. В данном варианте перенос отдельного виртуального окружения не уменьшил размер. Многоэтапная сборка позволяет исключать промежуточные инструменты и файлы, но сама по себе не гарантирует уменьшения образа.
\reportscreenshot{images/step-12.png}{Сравнение многоэтапной и предыдущих сборок}{fig:s12}

\chapter{Файл .dockerignore}
Создан .dockerignore с исключениями кеша Python, Git, виртуальных окружений, .env, логов и служебных файлов.
\reportscreenshot{images/step-13.png}{Содержимое .dockerignore}{fig:s13}

Его полный текст:
\begin{code}
__pycache__/
*.pyc
*.pyo
.git
.gitignore
.venv/
venv/
.env
*.log
Dockerfile*
README.md
\end{code}
После добавления исключений базовый Dockerfile пересобран с тегом v3.
\reportscreenshot{images/step-14.png}{Пересборка базового образа с тегом v3}{fig:s14}

Без .dockerignore в контекст сборки могли попасть .git, .env, локальные виртуальные окружения, логи и скомпилированные файлы Python. В образ они попадут только при соответствующих COPY или ADD. Здесь из корня копируется только requirements.txt, затем app/, поэтому служебные файлы корня не копировались и раньше. Размер v3 совпал с v1 --- 313 МБ.
Каталог .git хранит историю изменений: в ней могут оставаться ранее удалённые секреты и внутренние сведения. Включение истории в распространяемый образ способно раскрыть эти данные.

\chapter{Проверка здоровья контейнера}
\section{Настройка и сборка}
В итоговую стадию Dockerfile.optimized добавлена установка curl без рекомендуемых пакетов. Списки apt удаляются в том же слое.
Проверка HEALTHCHECK обращается к /health. Интервал проверок --- 30 секунд, таймаут --- 5 секунд, стартовый период --- 10 секунд, число неудачных проверок до unhealthy --- 3.
\reportscreenshot{images/step-15.png}{Dockerfile.optimized с curl и HEALTHCHECK}{fig:s15}

Финальный образ собран командой:
\begin{code}
docker build -t ivanovsky-assistant:final -f Dockerfile.optimized .
\end{code}
\reportscreenshot{images/step-16.png}{Успешная сборка финального образа}{fig:s16}

\section{Запуск и результат проверки}
Первый контейнер остановлен для освобождения порта 5451. Запущен ivanovsky-assistant-final с теми же персональными переменными окружения.
\reportscreenshot{images/step-19.png}{Остановка первого и запуск финального контейнера}{fig:s19}

После ожидания команда docker ps показала статус healthy, то есть встроенная проверка /health выполнилась успешно.
\reportscreenshot{images/step-17.png}{Статус healthy финального контейнера}{fig:s17}

Размер итогового образа составил 327 МБ против 307 МБ у ms. Добавление curl и его зависимостей увеличило отображаемый размер примерно на 20 МБ. Сравнение выполнено с ms, поскольку оба варианта используют одинаковую многоэтапную структуру.
\reportscreenshot{images/step-18.png}{Размеры образов после добавления curl}{fig:s18}

Без curl можно было использовать стандартный модуль Python urllib.request: он уже доступен в базовом Python-образе. Проверка должна завершаться ненулевым кодом при ошибке HTTP или соединения.

\chapter{Финальный артефакт и тегирование}
Созданы семантический тег 1.0.0 и тег даты 20261027.
\reportscreenshot{images/step-21.png}{Создание тегов версии и даты}{fig:s21}

Финальный артефакт содержит время выполнения 27.10.2026 18:01:37, логин пользователя, имя хоста, имя образа и таблицу тегов с размерами. Теги final, 1.0.0 и 20261027 имеют одинаковый идентификатор b6069c9bafc3: это ссылки на один и тот же образ, а не три независимые сборки.
\reportscreenshot{images/step-20.png}{Финальный артефакт выполнения}{fig:s20}

\reportelement{Заключение}
Собраны базовый образ v1, оптимизированный v2, вариант с многоэтапной сборкой ms, базовый образ v3 с .dockerignore и итоговый final с проверкой здоровья. Проверены ответы приложения, создание и чтение заметки, персонализация и секретный флаг. Финальный контейнер получил статус healthy.
\begin{table}[H]
\caption{Сводные результаты сборок}
\begin{tabularx}{\textwidth}{|l|X|r|}
\hline
Тег & Особенность & Размер, МБ \\\hline
v1 & Базовая установка зависимостей & 313 \\\hline
v2 & Без кеша pip & 287 \\\hline
ms & Multi-stage с виртуальным окружением & 307 \\\hline
v3 & Базовый Dockerfile с .dockerignore & 313 \\\hline
final & Multi-stage, curl и HEALTHCHECK & 327 \\\hline
\end{tabularx}
\end{table}
Все размеры приведены по округлённому столбцу DISK USAGE из Docker Desktop. Наилучший результат по размеру в проведённом сравнении показал v2. Итоговый образ дополнительно содержит curl и проверку здоровья.

\noindent Репозиторий работы на GitHub: \url{\RepositoryURL}.

\startreportappendices
\reportappendix{Финальный Dockerfile.optimized}
Короткие комментарии добавлены после выполнения сборок; команды и поведение Dockerfile не менялись.
\VerbatimInput[fontsize=\small,breaklines=true,breakanywhere=true]{Dockerfile.optimized}

\end{document}
